\section{Related Work}
\label{sec:relatedwork}

\subsection{Multimodal Embedding Models}

Vision-language contrastive models such as CLIP and ALIGN first showed that large-scale paired supervision can induce strong cross-modal alignment~\cite{CLIP,BLIP,blip2,SigLIP,SigLIP2,ALIGN}. Building on stronger MLLM backbones~\cite{Qwen25-VL,Qwen3VL,Qwen35Blog}, a growing line of work adapts generative MLLMs into unified multimodal embedders, including E5-V, GME, mmE5, VLM2Vec-V2, U-MARVEL, PDF-VLM2Vec and Qwen3-VL-Embedding~\cite{E5-V,GME,mme5,u-marvel,VLM2VecV2, pdfvlm2vec,Qwen3VLEmbedding}. These models extend retrieval from image--text pairs to videos, visual documents, and mixed-modality inputs, and are typically trained with contrastive learning over positive and negative query--document pairs~\cite{infonce,E5,GTE}, which is efficient, offline-encodable, and well aligned with industrial vector search. 
A parallel line of industrial systems scales this paradigm to production search
and recommendation, such as Pailitao-VL and
SAIL-Embedding~\cite{PailitaoVL,SAILEmbedding}, and the MOON series, which
progressively targets large-scale training~\cite{MOON}, modality
balance~\cite{MOON2}, and reasoning ability~\cite{MOON3}. These efforts confirm
the practical value of multimodal embedding under billion-scale traffic, yet
they remain centered on contrastive pair-level supervision.
However, such pair-level supervision only specifies which instances should be close, without modeling why they are relevant or which fine-grained evidence an embedding should preserve. DME follows the same efficient bi-encoder paradigm but supplements this pair-level objective with evidence-grounded and generative supervision to improve semantic sufficiency.

\subsection{Reasoning-Augmented Retrieval}
Chain-of-thought prompting and reinforcement learning have substantially
strengthened reasoning in multimodal large language models, through explicit
intermediate rationales, structured step-by-step solving, and reasoning-oriented
reward optimization~\cite{MMCoT,LLaVaCoT,VisionR1,R1Onevision,VCSTaR}. This
progress motivates a natural question for retrieval: whether such reasoning
ability can be transferred into the representations used for matching, rather
than kept only in free-form text generation.
A recent line of work reintroduces reasoning into retrieval representations. Think-Then-Embed generates embedding-centric reasoning traces before extracting a representation, and TRACE adaptively decides whether to generate a reasoning chain before emitting an embedding token~\cite{ThinkThenEmbed,TRACE}. UME-R1 explores reasoning-driven generative embeddings with retrieval-oriented reinforcement learning~\cite{UMER1}, while Embed-RL uses a frozen embedder as a reward model to train a reasoner that produces retrieval-useful chains of thought~\cite{EmbedRL}. These methods confirm that reasoning improves fine-grained retrieval, but most of them rely on explicit textual reasoning, additional generation, agentic tool use, or reranking-style computation, which increase query-encoding latency and complicate corpus indexing in billion-scale online systems. 
In contrast, DME performs reasoning entirely in the hidden space through a small number of retrieval-specific latent tokens, keeping inference close to a standard dense retriever while still acquiring evidence-aware reasoning ability.

\subsection{Generative and Reconstruction-based Representation Learning}
Reconstruction and generative objectives have long been used to strengthen representations, from masked language and image modeling to contrastive-captioning objectives~\cite{bert,MAE,CoCa}. In the multimodal embedding setting, recent studies show that content reconstruction or joint generative-retrieval training can push MLLMs to compress richer semantic information into embedding tokens~\cite{CoCoA,CREM}, and multi-token prediction has been shown to encourage representations that capture longer-range semantics~\cite{MultiTokenPrediction,DeepSeekV3,deepseekv4}. DME builds on this direction but uses the retrieval embedding itself as a semantic bottleneck for cross-conditional decoding, reconstructing counterpart-side content from the query and document embeddings. This not only injects fine-grained token-level supervision during training, but also makes the embedding decodable back to its input, which we use as a quantitative measure of semantic sufficiency rather than as an inference-time generation step.
